% RESEARCH_DRAFT_SHA256: 6547014cf673105c7092dd1288d935940b7a4150c54033dc34a9c9532cc5b97c
% Layout and source-reference wrappers expanded for the web renderer.
\documentclass[11pt]{article}
\usepackage[margin=25mm]{geometry}
\usepackage{fontspec,unicode-math}
\setmainfont{Latin Modern Roman}
\setmathfont{Latin Modern Math}
\usepackage{amsmath,amssymb,mathtools,xurl,hyperref}
\hypersetup{colorlinks=true,urlcolor=blue,linkcolor=blue}
\setcounter{secnumdepth}{0}
\setlength{\parindent}{0pt}
\setlength{\parskip}{5pt}
\setlength{\emergencystretch}{3em}
\newcommand{\sref}[2]{\hyperlink{source:#1:#2}{[S#2]}}
\newcommand{\cataloguescope}[1]{\paragraph{Source record.}#1}
\newcommand{\floor}[1]{\left\lfloor#1\right\rfloor}
\begin{document}
\section*{Research-attempt provenance}
This standalone research attempt retains the complete original definition body below,
including its historical source and preparation notes. Those notes describe the
earlier catalogue preparation; the new attempt and its verification are identified
separately after that body. No original ID, catalogue status, or source statement is changed.

\section{Duplicate Dimension-Free Separation-Exponent Problem}\label{statement}

\subsection{Definitions and mathematical statement}

\paragraph{Shared notation.}Standard number systems are $\mathbb N,\mathbb Z,\mathbb Q,\mathbb R,\mathbb C$. Specialized notation and index ranges follow the cited source and the definitions below.


A positive-semidefinite matrix has nonnegative quadratic form on every real vector. An affine set is a translate of a linear subspace. A separation exponent needs a specified inequality measuring the distance or rate of separation of these sets.
\paragraph{Record scope and qualifications.}
The catalogue title is retained as metadata. Named certificates, contracts, bounds or frameworks in that title are not asserted as results of the original source extract.
The supplied registry flags this extraction as unrecoverable. It is displayed as a literal source fragment; the definitions provide context rather than a replacement formal statement.

The relevant inequality and later mathematical statement are incomplete; the two joined excerpts are not treated as one formal theorem.

Source-record limitation: Source flags the statement as unrecoverable: The official report confirms that this open-problem-session extraction repeats the dimension-free separation-exponent problem represented canonically by OWR-14299088-010. No separate clean statement is promoted; the original record remains publicly visible and retained for provenance.
The source text below is rendered literally where necessary; damaged formulas and composite extracts are not silently repaired into a new mathematical claim.
The supplied catalogue marks this record as unrecoverable. Its extract and terminology are documented, but a complete original problem has not been reconstructed.
\paragraph{Statement recorded in the supplied source.}
Is there any quantity in algebraic geometry that can be used to express a dimension-free exponent for separation of an affine set and the set of positive semidefinite matrices? Is there a symbolic or a randomized algorithm to give such an exponent? Sum-of-Squares certificates for non-negativity of Quadratics Contributed by Grigoriy Blekherman Let A1 and A2 be two n \textbackslash{}times n symmetric real matrices, and let Q1 and Q2 be the corresponding quadratic forms Qi (x) = xt Ai x. It is well known that Q1 is nonnegative on the variety defined V defined by vanishing of Q2 if and only if there exists a scalar \textbackslash{}lambda such that the matrix A1 + \textbackslash{}lambdaA2 is positive semidefinite. This was rediscovered several times, and sometimes goes under the name of S-lemma. What happens if we instead ask for certificate of nonnegativity of Q1 on the variety defined by two quadratics Q2 and Q3 ? If there exist two scalar \textbackslash{}lambda2 , \textbackslash{}lambda3 such that A1 +\textbackslash{}lambda2 A2 +\textbackslash{}lambda3 A3 is positive semidefinite, then this will certify nonnegativity of Q1 , but existence of \textbackslash{}lambdai is no longer guaranteed. Instead we can look for higher degree certificates: find a sum of squares homogeneous polynomial P and homogeneous polynomials R2 and R3 all of degree 2d such that P Q1 + R2 Q2 + R3 Q3 is a sum of squares. The previously considered certificate had degree 2d equal to 0, and now we ask in what degree 2d is this certificate guaranteed to exist? Is the degree 2d of the certificate independent of the number of variables Optimization, Theoretical Computer Science and Algebraic Geometry 447 \href{https://doi.org/10.4171/owr/2025/9}{[S2]}
\paragraph{Source status.}
The supplied catalogue labels this source record open; that metadata was checked against the supplied snapshot on 2026-10-11. The label is preserved for provenance and is not an independent verification of present open status or mathematical correctness. The catalogue literature assessment is dated 2026-08-22.
\subsection{Short English statement}
The fragment asks for a dimension-independent separation exponent and then begins a quadratic-certificate discussion.
\subsection{Sources}
\begin{itemize}
\item[S1] Ulam AI and the UnsolvedMath Contributors, UnsolvedMath: A Curated Collection of Open Mathematics Problems, catalogue entry. Catalogue curation and metadata: CC BY 4.0. The source extract is retained; definitions, English summaries and qualifications are editorial additions.
\url{https://huggingface.co/datasets/ulamai/UnsolvedMath}.
\item[S2] MFO-RIMS Tandem Workshop: Optimization, Theoretical Computer Science and Algebraic Geometry: Convexity and Beyond (2025), Oberwolfach Reports, DOI 10.4171/owr/2025/9, pp. 405--450
\url{https://doi.org/10.4171/owr/2025/9}.
\item[S3] Ochiai--Sekiguchi--Waki
\url{https://doi.org/10.1016/j.jmaa.2024.129070}.
\item Source-specific Oberwolfach Reports licences, including CC BY-SA 4.0 where applicable, are retained; see THIRD\_PARTY\_NOTICES.md.
\end{itemize}
\paragraph{Preparation and review.}
The original source and UnsolvedMath contributors retain attribution for the source extract. Codex (GPT-6; exact serving revision not exposed) assisted with definitions, English summaries, scope notes and formatting. No new human mathematical authorship or independently verified solution is asserted. Catalogue identity checked against the complete supplied registry and current repository. Distinct upstream records, including repeated questions and broad or damaged extracts, are retained for complete coverage of the source CS category. A repeated record is not a new independent conjecture. Complete mathematical-equivalence and current-literature review has not been performed. Preparation and final review follow contribution prompts C and D, with source gaps disclosed. Structural and rendering checks do not establish a new theorem or open status.
\label{end:statement}

\section{Literature and context check}
The primary report, printed pp. 446--447, separates the quadratic-certificate discussion from the affine/PSD separation question and gives the missing distance inequality and a singularity-degree-dependent exponent. The retained catalogue record is a damaged duplicate joining those discussions. The construction below concerns a precise local form of the separation question only; it is not a reconstruction of a new independent problem from the joined fragment.
\section{Research attempt: An explicit chain forcing large local separation exponents}

For an integer $m\geq1$, consider symmetric block-diagonal $2m\times2m$ matrices $A(z)$ with variables $z_1,\ldots,z_m$ and fixed $z_0=0$, whose $i$th block, $0\leq i<m$, is
\[
B_i(z)=\begin{pmatrix}z_i&z_{i+1}\\z_{i+1}&1\end{pmatrix}.
\]
Their set $E$ is affine. A block is PSD if and only if $z_i\geq z_{i+1}^2$, as follows by completing the square in $z_i u^2+2z_{i+1}uv+v^2$. Since $z_0=0$, positivity of the first block forces $z_1=0$, then positivity of the next forces $z_2=0$, and so on. Thus $E\cap S_+^{2m}$ is the singleton $A(0)$.

For $0<t\leq1$, set
\[
z_i(t)=t^{2^{m-i}},\qquad 1\leq i\leq m.
\]
Every block except the first is PSD of determinant zero, because $z_i=z_{i+1}^2$ for $i\geq1$. The first block becomes PSD if its upper-left entry is changed from zero to $z_1^2=t^{2^m}$. This is a correction of Frobenius norm exactly $t^{2^m}$. Consequently
\[
\operatorname{dist}(A(z(t)),E)=0,\qquad
\operatorname{dist}(A(z(t)),S_+^{2m})\leq t^{2^m}.
\]
On the other hand the last off-diagonal entries are $z_m=t$, so
\[
\operatorname{dist}(A(z(t)),E\cap S_+^{2m})
=\|A(z(t))-A(0)\|_F\geq\sqrt2\,t.
\]
The family is bounded for fixed $m$. Therefore any local inequality
\[
\operatorname{dist}(Z,E)+\operatorname{dist}(Z,S_+^{2m})
\geq c\,\operatorname{dist}(Z,E\cap S_+^{2m})^\gamma,
\qquad c>0,
\]
in a neighborhood of $A(0)$ must have $\gamma\geq2^m$: otherwise substituting the curve and dividing by $t^\gamma$ gives a positive constant bounded above by $t^{2^m-\gamma}\to0$. The denominator in the primary report's global inequality, $1+\|Z\|^2$, stays bounded above and below along this curve, so it does not change this necessary exponent comparison.

This rules out an absolute exponent working for all matrix dimensions and all affine sections. It does not rule out the source's sought expression in geometric quantities, which may themselves grow with dimension or singularity depth; indeed the source already discusses such dependence. The displayed chain explicitly exposes the repeated quadratic tangency that a proposed invariant or algorithm would have to detect.

\paragraph{The joined quadratic fragment.}
It remains a separate statement. For instance, if $Q_2\ne0$, taking $P=Q_2^2$, $R_2=-Q_1Q_2$, $R_3=0$ makes its displayed certificate identically zero regardless of the sign of $Q_1$. Thus the omitted multiplier convention matters and cannot be inferred from the separation calculation.
\section{Verification and adversarial audit}
The entire affine/PSD intersection was computed, not just its intersection with the chosen curve. The PSD upper-distance bound uses an explicit single-entry correction, so no unstated projection formula is required. $m=1$ reduces to the familiar quadratic tangency example. The exponent argument is local and accommodates any fixed positive constant and any bounded normalization denominator.
\paragraph{Audit corrections.} Distinguished an absolute dimension-independent exponent from an exponent described using dimension-dependent geometric data; preserved the duplicate/damaged status of the original extract.
\paragraph{External computation.} None; symbolic arguments and boundary checks only.
\section{Scope and remaining gaps}
No dimension-free algebraic expression or symbolic/randomized algorithm for the intended source exponent is provided. The damaged duplicate and the adjoining certificate text remain visibly separate; the auxiliary lower-bound family does not receive full-target completion.
\paragraph{Integration review.} Codex (GPT-6; exact serving revision not exposed) reviewed the written arguments and expanded layout and source-reference wrappers. Repository metadata and the production web renderer were checked using Python 3.13 and JavaScript with MathJax 3.2.2. This model-assisted review is not independent human proof verification; no PDF compilation or formal proof checking is asserted.

\section{Completion Estimate}
\noindent\textbf{COMPLETION: 0\%}
The calculations below the preserved source statement investigate a concrete route,
but establish no substantiated new progress toward the entire original target.
Known special cases and failed sufficient conditions are not credited as a solution.
\section{References}
The original source references and catalogue attribution are preserved in the
Sources section above. The following URLs were actually consulted for this attempt;
retention of an additional original reference does not claim fresh consultation of it.
\begin{itemize}
\item \url{https://ems.press/content/serial-article-files/51353} (consulted 2026-10-11). 
\end{itemize}
\section{Model and contributor disclosure}
Codex GPT-6 completed written-proof review of this attempt at its declared scope; this is model-assisted review, not independent human proof checking.
\paragraph{Material model collaboration.} Root Codex GPT-6 supplied a corroborating audit of the vacuous nonzero-multiplier identity in the separate quadratic fragment; the repeated-tangency separation construction and proof were developed by the Astra author. The gpt-6-astra author inspected and reworked these arguments; this is model-assisted mathematical work, not independent human verification.
OpenAI gpt-6-astra, requested reasoning configuration ultra, generated this research
attempt on 2026-10-11. The exact serving revision is not exposed. Codex GPT-6
performed orchestration and repository-format review; its exact serving revision
is also not exposed. The mathematical route, boundary checks and final adversarial
self-audit followed contribution prompts A then B in
\texttt{docs/contribution\_prompts.md}, after reading \texttt{CONTRIBUTING.md}.
Python 3.14.2 assembled the standalone documents; any mathematical computation is
specified in Verification. No theorem prover was used. No human mathematical
authorship, substantial human proof correction, or independent human proof checking
is asserted. Model-assisted review, source retrieval, and structural checks do not
establish correctness or novelty. Full private reasoning transcripts are unavailable
as reproducibility artifacts; the complete arguments, inputs and outcomes actually
used are stated here or in the companion batch audit.
\section{Final status}
Unresolved as a new research contribution. This is a concrete route investigation,
not a claim of new complete resolution or a current-status certification. A known
answer reproduced for context does not become a new contribution. The catalogue's
source-status label is preserved; this attempt status does not assert that a
source-labelled solved question is open.
\end{document}
